
Large language models (LLMs) are increasingly deployed as agentic  systems that reason over context \cite{giurgiu2025supporting}, plan multi-step actions, and invoke external tools \cite{wang2025agent} to accomplish user goals.The Model Context Protocol (MCP) has emerged as a common substrate for such tool-augmented agents by standardizing how LLM applications connect to external tools and data sources \cite{anthropic_mcp_2024}.
MCP introduces a host--client--server architecture in which tool providers expose capabilities via MCP servers, and agent hosts dynamically discover and invoke those capabilities through standardized messages, schemas, and tool descriptors \cite{mcp2025spec}. This design addresses integration fragmentation and encourages a rapidly growing ecosystem of third-party servers and reusable connectors \cite{anthropic_mcp_2024}. At the same time, MCP shifts the attack surface from a single prompt boundary to a distributed, multi-party context and tool supply chain \cite{errico2025securing}.



The rapidly emerging research area of LLM agent security has recently shifted focus toward the security of MCP \cite{jing2025mcip}. While MCP significantly enhances agent capabilities, the common motivation across the reviewed literature is the critical realization that this architecture introduces severe, novel attack vectors, such as server-side tool poisoning, host-side prompt injection, and unauthorized privilege escalation \cite{zong2025mcp}, because models inherently trust the context and tools provided by third-party MCP servers. To address this problem, these papers collectively pioneered the techniques of systematic threat modeling and empirical vulnerability evaluation, introducing the first comprehensive benchmarking datasets (such as MSB\cite{zhang2025msb}, MCP-SafetyBench \cite{zong2025mcp}, and MCPSecBench \cite{yang2025mcpsecbench}) alongside novel defensive frameworks like the Model Contextual Integrity Protocol (MCIP) \cite{jing2025mcip}. These specific papers were selected for inclusion because they represent the foundational, state-of-the-art research defining this brand-new domain; together, they establish the taxonomy of MCP threats, provide the empirical ground truth for evaluating current LLM robustness against these exploits, and propose the initial generation of mitigation strategies.
